%==============================================================
% ABSTRACT (Περίληψη)
%==============================================================
\chapter*{Περίληψη}
\addcontentsline{toc}{chapter}{Περίληψη}

Η ραγδαία υιοθέτηση μοντέλων μηχανικής μάθησης στον τομέα του
πιστωτικού κινδύνου εισάγει ένταση μεταξύ της αυξημένης προβλεπτικής
ικανότητας και της μειωμένης ερμηνευσιμότητας, σε ένα κανονιστικό
πλαίσιο που γίνεται διαρκώς πιο απαιτητικό. Η παρούσα εργασία
εξετάζει πώς επικυρώνονται, ερμηνεύονται και παρακολουθούνται τα
σύγχρονα μοντέλα μηχανικής μάθησης στον πιστωτικό κίνδυνο και πώς
μεταβάλλονται οι απαιτήσεις αυτές καθώς αυξάνεται η πολυπλοκότητα
του μοντέλου. Για τον σκοπό αυτό συγκρίνονται τρία μοντέλα
διαφορετικής πολυπλοκότητας στη βάση δεδομένων FICO HELOC: η
Logistic Regression ως παραδοσιακό σημείο αναφοράς, το XGBoost ως
κύριο μοντέλο ανάλυσης και ένα Feed-Forward Neural Network ως
benchmark αυξημένης πολυπλοκότητας. Η αξιολόγηση πραγματοποιείται
σε τέσσερις άξονες: απόδοση, ερμηνευσιμότητα, ανίχνευση μεροληψίας
και παρακολούθηση σταθερότητας. Τα μοντέλα παρουσιάζουν οριακές
διαφορές στους δείκτες απόδοσης (AUC 0,789--0,796), εύρημα που
εξηγείται από την εκτεταμένη γραμμικότητα της βάσης λόγω
προεπεξεργασίας από την FICO. Η ανίχνευση μεροληψίας με τη
βιβλιοθήκη Fairlearn αποκαλύπτει σχεδόν ταυτόσημες τιμές DPD
(${\sim}0{,}33$) και DI (${\sim}0{,}47$) και στα τρία μοντέλα,
αποδεικνύοντας ότι η πηγή της μεροληψίας είναι τα δεδομένα και όχι
η εσωτερική λογική κάποιου αλγορίθμου. Ο μετριασμός μεροληψίας
συνεπάγεται κόστος που διαφοροποιείται ανά τεχνική: η
in-processing προσέγγιση μειώνει την AUC της Logistic Regression
κατά 14\%, ενώ η post-processing προσέγγιση διατηρεί την AUC του
XGBoost αμετάβλητη με μικρή απώλεια στο Balanced Accuracy. Η
παρακολούθηση μέσω PSI και CSI σε προσομοιωμένα σενάρια οικονομικής
επιδείνωσης λειτουργεί ως αξιόπιστο σύστημα έγκαιρης προειδοποίησης.
Ωστόσο, η ερμηνευσιμότητα παρουσιάζει τη δραματικότερη διαφοροποίηση
ανά μοντέλο: η ανάλυση SHAP για το XGBoost ολοκληρώνεται σε
δευτερόλεπτα μέσω TreeExplainer, ενώ για το FFNN απαιτεί
${\sim}2{,}5$ ώρες μέσω KernelExplainer. Με βάση τα ευρήματα
προτείνεται πλαίσιο επικύρωσης που περιλαμβάνει αρχές σχεδιασμού,
πρότυπη ροή εργασιών επτά βημάτων και διαφοροποιημένες απαιτήσεις
ανά επίπεδο πολυπλοκότητας. Η κεντρική συμβολή του πλαισίου είναι η
απόδειξη ότι η αύξηση της πολυπλοκότητας δεν αυξάνει γραμμικά όλες
τις απαιτήσεις επικύρωσης --- η ερμηνευσιμότητα είναι ο άξονας στον
οποίο η πολυπλοκότητα έχει δραματική επίπτωση, ενώ η παρακολούθηση
και η ανίχνευση μεροληψίας παραμένουν σχετικά σταθερές.

\vspace{2cm}

\noindent\textbf{Λέξεις-κλειδιά:} πιστωτικός κίνδυνος, μηχανική
μάθηση, επικύρωση μοντέλων, ερμηνευσιμότητα, ανίχνευση μεροληψίας,
παρακολούθηση μοντέλων, SHAP, XGBoost, νευρωνικά δίκτυα,
FICO HELOC, Basel II/III.

%--------------------------------------------------------------
% ENGLISH ABSTRACT
%--------------------------------------------------------------
\chapter*{Abstract}
\addcontentsline{toc}{chapter}{Abstract}
\begin{english}

The rapid adoption of machine learning models in credit risk
introduces a tension between increased predictive performance and
reduced interpretability, within a regulatory framework that is
becoming progressively more demanding. This thesis examines how
modern machine learning models are validated, interpreted and
monitored in credit risk, and how these requirements scale as
model complexity increases. To this end, three models of varying
complexity are compared on the FICO HELOC dataset: Logistic
Regression as a traditional baseline, XGBoost as the primary
analytical model, and a Feed-Forward Neural Network as a
high-complexity benchmark. Evaluation is carried out across four
axes: performance, interpretability, bias detection and stability
monitoring. The models exhibit marginal differences in performance
metrics (AUC 0.789--0.796), a finding explained by the extensive
linearity of the dataset due to FICO's preprocessing. Bias
detection using the Fairlearn library reveals nearly identical
DPD ($\sim$0.33) and DI ($\sim$0.47) values across all three models,
demonstrating that the source of bias lies in the data rather than
in any model's internal logic. Bias mitigation comes at a cost
that varies by technique: the in-processing approach reduces the
AUC of Logistic Regression by 14\%, while the post-processing
approach leaves the AUC of XGBoost unchanged at the cost of a
modest reduction in Balanced Accuracy. Monitoring through PSI and
CSI on simulated economic deterioration scenarios proves to be a
reliable early warning system. However, interpretability shows the
most dramatic differentiation across models: SHAP analysis for
XGBoost completes in seconds via TreeExplainer, whereas for the
FFNN it requires approximately 2.5 hours via KernelExplainer.
Based on these findings, a validation framework is proposed
comprising design principles, a seven-step standard workflow and
differentiated requirements per complexity level. The central
contribution of the framework is the demonstration that increased
model complexity does not scale all validation requirements
linearly --- interpretability is the axis on which complexity has
dramatic impact, while monitoring and bias detection remain
relatively stable.

\vspace{2cm}

\noindent\textbf{Keywords:} credit risk, machine learning, model
validation, interpretability, bias detection, model monitoring,
SHAP, XGBoost, neural networks, FICO HELOC, Basel II/III.

\end{english}